\documentclass[11pt]{article}
\usepackage[margin=1in]{geometry}
\usepackage{amsmath,amssymb,amsthm}
\usepackage{graphicx}
\usepackage[hidelinks]{hyperref}
\setlength{\emergencystretch}{2em}

\newtheorem{theorem}{Theorem}
\newtheorem{corollary}{Corollary}
\newtheorem{definition}{Definition}
\newtheorem{remark}{Remark}

\newcommand{\C}{\mathbb C}
\newcommand{\dist}{\operatorname{dist}}
\newcommand{\otmin}{\otimes_{\min}}

\title{A Relative Two-Evaluation Criterion for Countable Capture}
\author{Run \texttt{fa\_banach\_001}, lane 19}
\date{June 17, 2026}

\begin{document}
\maketitle

\section*{Source Problem}

The source paper is Jinghao Huang, Karimbergen Kudaybergenov, and Rui Liu,
\emph{The ball-covering property of von Neumann algebras and noncommutative
symmetric spaces}, arXiv:2311.14257.  After proving the commutative
\(C(K)\) criterion and the von Neumann algebra criterion, it asks for a
noncommutative \(C^*\)-algebra version:

\begin{center}
\includegraphics[width=0.92\textwidth]{figures/open_problem_crop.png}
\end{center}

\section*{Statement}

\begin{definition}
Let \(A\) be a unital \(C^*\)-algebra.  We say that \(A\) has the
\emph{relative two-evaluation property} if for every separable unital
\(C^*\)-subalgebra \(B\subseteq A\), there are a self-adjoint element
\(u\in A\) with \(\|u\|=1\), a unital \(C^*\)-algebra \(D\), and contractive
linear maps
\[
        \theta_+,\theta_-: C^*(B,u)\to D
\]
such that the restrictions \(\theta_\pm|_B\) are isometric and agree, while
\[
        \theta_+(u)=1_D,\qquad \theta_-(u)=-1_D.
\]
\end{definition}

\begin{theorem}
\label{thm:main}
If a unital \(C^*\)-algebra \(A\) has the relative two-evaluation property,
then \(A\) has countable capture.  Consequently \(A\) fails BCP.
\end{theorem}

\begin{corollary}[tensorial sufficient condition]
\label{cor:tensorial}
Suppose that for every separable unital \(C^*\)-subalgebra \(B\subseteq A\)
there is a self-adjoint \(u\in A\) with \(\|u\|=1\), \(-1,1\in\sigma(u)\),
\(u\in B'\cap A\), and the canonical multiplication map
\[
        B\otmin C^*(u)\to C^*(B,u)
\]
is an isomorphism.  Then \(A\) has countable capture and fails BCP.
\end{corollary}

\section*{Idea of the Proof}

The countable-capture characterization says that it is enough, after seeing a
countable centre set \(C\), to find \(u\in S_A\) such that
\[
        \dist(c,\C u)=\|c\|\qquad(c\in C).
\]
Put the countable set inside a separable subalgebra \(B\).  The two evaluation
maps turn \(c-\lambda u\) into two elements
\[
        \theta(c)-\lambda 1,
        \qquad
        \theta(c)+\lambda 1.
\]
The maximum of their norms is at least \(\|\theta(c)\|=\|c\|\), because their
sum is \(2\theta(c)\).  This forces
\[
        \|c-\lambda u\|\ge \|c\|
\]
for every scalar \(\lambda\), hence gives countable capture.

\section*{Proof}

\begin{proof}[Proof of Theorem~\ref{thm:main}]
Let \(C\subset A\) be countable.  Let \(B=C^*(1,C)\), a separable unital
\(C^*\)-subalgebra of \(A\).  Choose \(u,D,\theta_+,\theta_-\) from the
relative two-evaluation property.  Since \(\theta_\pm|_B\) agree, write
\(\theta:B\to D\) for their common restriction.  This restriction is
isometric.

Fix \(c\in C\) and \(\lambda\in\C\).  Contractivity gives
\[
\begin{aligned}
        \|c-\lambda u\|
        &\ge
        \max\{\|\theta_+(c-\lambda u)\|,
              \|\theta_-(c-\lambda u)\|\}\\
        &=
        \max\{\|\theta(c)-\lambda1\|,
              \|\theta(c)+\lambda1\|\}.
\end{aligned}
\]
By the triangle inequality,
\[
        2\|c\|
        =
        2\|\theta(c)\|
        =
        \|(\theta(c)-\lambda1)+(\theta(c)+\lambda1)\|
        \le
        2\max\{\|\theta(c)-\lambda1\|,\|\theta(c)+\lambda1\|\}.
\]
Thus \(\|c-\lambda u\|\ge\|c\|\) for every \(\lambda\).  Taking the infimum
over \(\lambda\) gives \(\dist(c,\C u)\ge\|c\|\), and equality follows from
\(\lambda=0\).  Therefore \(A\) has countable capture by the previous
countable-capture characterization.  The countable-capture obstruction then
implies that \(A\) fails BCP.
\end{proof}

\begin{proof}[Proof of Corollary~\ref{cor:tensorial}]
Let \(B\subseteq A\) be separable and choose \(u\) as in the hypothesis.
Because \(-1,1\in\sigma(u)\), the evaluation maps
\[
        \operatorname{ev}_{\pm1}:C^*(u)\cong C(\sigma(u))\to\C
\]
are characters.  Under the isomorphism
\[
        C^*(B,u)\cong B\otmin C^*(u),
\]
the maps
\[
        \theta_\pm=\operatorname{id}_B\otimes\operatorname{ev}_{\pm1}
\]
are contractive and fix \(B\), while sending \(u\) to \(\pm1\).  Hence the
relative two-evaluation property holds, and Theorem~\ref{thm:main} applies.
\end{proof}

\section*{Scope}

This criterion recovers the fresh-coordinate argument for uncountable tensor
products.  More importantly, it isolates the exact kind of lemma that would
settle the remaining arbitrary simple nonseparable \(C^*\)-algebra wall: one
would need to show that every separable subalgebra admits such a fresh
two-evaluation direction, or else construct a counterexample where no such
direction exists.

\section*{Novelty and Search Bounds}

The local run indexes and web searches were checked for countable capture,
two-evaluation, infinite tensor products, relative commutants, and BCP for
\(C^*\)-algebras.  No exact statement of this relative criterion was found.
The theorem is a distilled abstraction of the uncountable tensor-product
capture packet.

\section*{Human Review Focus}

Reviewers should check:
\begin{enumerate}
    \item the definition of the two-evaluation property;
    \item the distance inequality in Theorem~\ref{thm:main};
    \item the tensorial sufficient condition in Corollary~\ref{cor:tensorial}.
\end{enumerate}

\begin{thebibliography}{9}

\bibitem{HKL}
Jinghao Huang, Karimbergen Kudaybergenov, and Rui Liu,
\emph{The ball-covering property of von Neumann algebras and noncommutative
symmetric spaces},
arXiv:2311.14257, 2023.

\bibitem{lane19capture}
Run \texttt{fa\_banach\_001}, lane 19,
\emph{When the Countable-Capture BCP Obstruction Exists},
solution packet
\texttt{2311.14257\_countable\_capture\_characterization}, 2026.

\bibitem{lane19tensor}
Run \texttt{fa\_banach\_001}, lane 19,
\emph{Uncountable Tensor Products and Countable Capture},
solution packet
\texttt{2311.14257\_uncountable\_tensor\_product\_capture}, 2026.

\end{thebibliography}

\end{document}
